
\section{The contraction lemma and the original two-step conclusion}
\label{sec:sr-source-two-step}

This section translates
\srcid{EigenvalueEndgame.SignatureSplitting.ConvCont},
\srcid{SpanPairIff}, and \srcid{TwoStepOf}.
The contraction is first proved for an abstract locally compact group.
Its application concludes that a lower restriction has a pure
\(\alpha\) constant term. Proportionality and nonvanishing of the
lower theta images are not hypotheses of this argument.

Let \(G\) be a second-countable locally compact Hausdorff topological
group, with its Borel structure and left Haar measure \(\mu\).
Let \(X\) be a topological measurable space whose open sets are
measurable, with probability measure \(\lambda\).
Fix a subgroup \(\Gamma\subset G\), a continuous map
\(n:X\to G\), and \(a\in G\).
For a neighborhood \(V\) of \(1\), call a real function \(\eta\)
a bump in \(V\) if it is continuous, compactly supported,
nonnegative, has support contained in \(V\), and \(\eta(1)\ne0\).
No normalization of its integral is imposed.

The source's cusp condition consists of the following two assertions.
There is a compact \(\Omega_0\subset G\) such that every \(g\in G\)
has \(\gamma g=a^k\kappa\), with \(\gamma\in\Gamma\),
\(k\ge0\), and \(\kappa\in\Omega_0\).
Moreover every neighborhood \(V\) admits a bump \(\eta\) such that,
for every compact \(\Omega\subset G\), there are a real \(D\)
and a compact \(\Omega'\subset G\) with this property:
for every \(x\in X\) and \(k\ge0\), some \(\gamma\in\Gamma\)
and \(u\in\Omega'\) satisfy
\[
 \gamma n(x)a^k=a^ku,\qquad
 |\eta(\kappa^{-1}u^{-1}y)-\eta(\kappa^{-1}y)|
                   \le D2^{-k}
 \quad(\kappa\in\Omega,\ y\in G).                 \tag{S1}
 \label{eq:sr-source-bump-cusp}
\]
For \(m\ge0\), say that \(\phi:G\to\mathbb C\) has growth
\(m\) along \(a\) if, for every compact \(\Omega\), some \(C\)
satisfies
\[
       |\phi(a^k\kappa)|\le C(2^k)^m
                 \quad(k\ge0,\ \kappa\in\Omega). \tag{S2}
 \label{eq:sr-source-growth-along-a}
\]

\begin{srlemma}[The abstract contraction, \srcid{ConvCont.kill}]
\label{lem:sr-source-abstract-contraction}
Assume the cusp condition. Let \(\phi,c:G\to\mathbb C\)
be continuous, let \(\phi\) be left \(\Gamma\)-invariant,
and suppose
\[
 c(g)=\int_X\phi(n(x)g)\,d\lambda(x),\qquad
 c(a^2g)-c(ag)=z\{c(ag)-c(g)\}                   \tag{S3}
 \label{eq:sr-source-contraction-data}
\]
for every \(g\), where \(|z|<1\).
If \(\phi\) has growth \(m\) for some integer \(m\ge0\),
then \(c(ag)=c(g)\) for every \(g\).
\end{srlemma}

\begin{proof}
Put \(Dc(g)=c(ag)-c(g)\).
Induction on \(k\) in (S3) gives
\[
 Dc(a^kg)=z^kDc(g),\qquad
 c(a^kg)=c(g)+(1+z+\cdots+z^{k-1})Dc(g).
 \tag{S4}\label{eq:sr-source-contraction-powers}
\]
Continuity on compact sets and the convergent geometric
majorant show that \(c\) has growth zero.

For a bump \(\eta\), define
\(\phi_\eta(g)=\int_G\phi(gy)\eta(y)\,d\mu(y)\),
and similarly \(c_\eta\).
These convolutions are continuous and left
\(\Gamma\)-invariant on the \(\phi\) side.
Their unipotent-average and recurrence identities are again (S3).
Here the needed Fubini justification does not assume
compactness of \(X\): (S1), with \(k=0\), puts every
\(n(x)\) modulo \(\Gamma\) into one compact set.
Thus \(\phi(n(x)g)\) is bounded uniformly in \(x\)
when \(g\) ranges over a compact set.
Multiply this bound by the compactly supported \(\eta\)
to obtain an integrable product majorant.

Suppose now that \(\phi\) has growth \(m+1\).
Choose \(\eta\) satisfying (S1) for an arbitrary prescribed
neighborhood \(V\).
For \(\kappa\) in a compact \(\Omega\), change variables
by left multiplication in the convolution and use
\(\Gamma\)-invariance. It gives
\[
 \phi_\eta(a^k\kappa)-\phi_\eta(n(x)a^k\kappa)
 =\int_G\phi(a^ky)
       \{\eta(\kappa^{-1}y)-\eta(\kappa^{-1}u^{-1}y)\}\,d\mu(y).
\]
Both terms are supported in the compact set
\[
          K=\Omega'\Omega\operatorname{supp}\eta
                          \ \cup\ \Omega\operatorname{supp}\eta.
\]
The growth \(m+1\) bound on \(K\) and (S1) bound this integral
by \(C D\mu(K)(2^k)^m\).
Averaging in \(x\), using the probability measure, expresses
\(\phi_\eta(a^k\kappa)\) as \(c_\eta(a^k\kappa)\) plus
an error with this bound. The recurrence shows that
\(c_\eta\) has growth zero, so \(\phi_\eta\) has growth \(m\).
This is precisely \srcid{Dat.step}.

We induct on \(m\).
For \(m=0\), the compact ray covering and automorphy make
\(\phi\) globally bounded by some \(M\).
Its probability average is bounded by the same \(M\).
Applying (S4) at \(a^{-k}g\) gives
\(|Dc(g)|\le2M|z|^k\), hence \(Dc=0\).

For the induction step, the growth drop and induction hypothesis
give \(Dc_\eta=0\) for a suitable bump in every neighborhood of \(1\).
The source removes the convolution by positivity.
If \(D_0(g)=Dc(g)\ne0\), continuity makes
\(\operatorname{Re}(D_0(gy)/D_0(g))>0\) in a neighborhood of \(1\).
Choose its nonnegative bump. Its value at \(1\) is positive,
and Haar measure has full support, so
\[
 \operatorname{Re}\int_G
           \frac{D_0(gy)}{D_0(g)}\eta(y)\,d\mu(y)>0.
\]
This contradicts \(Dc_\eta(g)=0\).
The proof therefore needs neither an integral-one bump
nor a previously established approximate-identity theorem.
\end{proof}

We now specify the exact application.
For a fixed triple let
\[
 \operatorname{Sig}(G)=
 \{I_V(\Phi)-I_{V'}(\Phi'):
       \Phi\in\operatorname{PolData}(M),\
       \Phi'\in\operatorname{PolData}(M')\cap\operatorname{So}(M')\}.
\]
No matching relation is imposed in this definition.
For a target hyperbolic plane \(H'\) and its complement \(W_1\),
write \(\operatorname{Split}_2(F;F_\alpha,F_\beta)\) for
its line constant term split into the laws
\(\alpha(t)|t|^{a'/2}\) and \(\beta(t)|t|^{a/2}\).
Write
\[
 \operatorname{IsLow}(H',F_\alpha,k,f_1)
 \quad\Longleftrightarrow\quad
 f_1(y)=F_\alpha(lk)
       \text{ whenever }l\text{ is a Levi lift of }(1,y).
\]
The set \(\operatorname{FinK}(W)\) consists of all adelic
target elements commuting with the prescribed Cartan
at every infinite place. It does not require their finite
components to lie in a particular compact subgroup.
Finally \(\operatorname{Pure}(W_1,\alpha,a'/2,f_1)\) means
that there exist a further nondegenerate rational target
plane, its complement \(W_2\), and a constant term of \(f_1\)
with the single law \(\alpha(t)|t|^{a'/2}\).

Two intermediate assertions are retained exactly.
\(\operatorname{SigMod}\) says that every member of
\(\operatorname{Sig}(G)\) has polynomial growth in the actual
adelic height. \(\operatorname{SigSplit}\) says that whenever
the signature sum is positive there is one nondegenerate
rational target plane, chosen before the form, along which
every member of \(\operatorname{Sig}(G)\) admits a two-component
splitting. These are not assertions of proportionality.

\begin{srproposition}[The original \srcid{TwoStep} conclusion]
\label{prop:sr-source-two-step}
Assume \(\operatorname{SigMod}\), \(\operatorname{SigSplit}\),
the target Iwasawa assertion \(\operatorname{Iwa}\), and
automorphy of both polynomial period images
\(\operatorname{IAut}\).
Let \(G\) have ranks \((b+2,a')\), and let \(G_1\) have
ranks \((b,a')\), with their specified triples, and assume explicitly
\(a'<b\). The triples' character data are unitary: in particular,
\(|\alpha(t)|=|\beta(t)|=1\) for every idèle \(t\).
Assume the exact plane-removal configuration:
the small first space is unchanged, \(q_1=q\),
\(p_1+1=p\), and the large first and target planes
are adapted in columns \(0,b+1\).
Suppose \(q\ge1\) and \(F\in\operatorname{Sig}(G)\) has
\[
             \operatorname{Split}_2(F;F_\alpha,0).
\]
Assume that for every \(k\in\operatorname{FinK}(W)\) some
\(f_1\in\operatorname{Sig}(G_1)\) satisfies
\(\operatorname{IsLow}(H',F_\alpha,k,f_1)\).
Then every such \(k,f_1\) satisfies
\[
             \operatorname{Pure}(W_1,\alpha,a'/2,f_1).
 \tag{S5}\label{eq:sr-source-two-step-pure}
\]
\end{srproposition}

\begin{proof}
Since \(q_1=q\ge1\), \(\operatorname{SigSplit}\) supplies one
further rational plane \(H_1\) in \(W_1\), with nondegenerate
complement \(W_2\), along which every lower comparison form
splits. Choose its two components \(A_x,C_x\) for the function
\(y\mapsto F_\alpha(\operatorname{lev}(y)x)\), for every upper \(x\).

Here is why this choice is possible for arbitrary \(x\).
The upper Iwasawa assertion gives
\(x=n\,m(t)\operatorname{lev}(g_1)k\) with
\(k\in\operatorname{FinK}(W)\).
The constant term \(F_\alpha\) is left invariant under
the line radical, and its outer Levi law gives the factor
\(\alpha(t)|t|^{a'/2}\).
Its restriction is consequently a scalar multiple of a
right translate of the supplied lower \(f_1\).
Transporting the lower splitting by that scalar and translation
produces \(A_x,C_x\).
Their laws are \(\alpha(t)|t|^{a'/2}\) and
\(\beta(t)|t|^{b/2}\).
The value \(c(x)=A_x(1)+C_x(1)\) is independent of this choice:
it is the lower constant-term average of the displayed
restriction. This uses uniqueness of the average, not
an unmentioned choice of individual component normalizations.

The two target planes give an embedding of
\(\mathrm{SL}_2(\mathbb A_E)\) acting on their isotropic
two-plane and by inverse conjugate transpose on its dual.
Put \(u=(\operatorname{halfE})^{-1}\), whose infinite components
are \(2\), and \(r=|u|_E>1\).
Let \(h\) be the image of \(\operatorname{diag}(u,u^{-1})\).
The two successive Levi laws give
\[
 c(hx)=A_x(1)+zC_x(1),\quad
 c(h^2x)=A_x(1)+z^2C_x(1),\quad
 z=\alpha(u)\beta(u^{-1})r^{(a'-b)/2}.
 \tag{S6}\label{eq:sr-source-two-step-recurrence}
\]
Since \(|\alpha(u)|=|\beta(u^{-1})|=1\) and \(a'<b\),
\(|z|=r^{(a'-b)/2}<1\).
This is exactly where the lower rank inequality is used.

The source's \srcid{SpanPairIff.two\_orders} identifies this
iterated line average with a different staged average.
First average \(F\) over the radical attached to the rational
isotropic two-plane. Then average its restriction to the
embedded \(\mathrm{SL}_2\) over its upper unipotent quotient.
The equality is proved by the two quotient maps of the
flag unipotent group: one has the first line radical as
kernel and the lower line radical as quotient; the other
has the two-plane radical as kernel and the additive
\(E\)-root group as quotient.
Their rational quotients are compact and their Haar measures
are normalized to mass one, so the two stage formulas have
the same normalization.

The resulting \(\mathrm{SL}_2\) function is continuous and
rationally invariant. Its polynomial ray growth follows
from \(\operatorname{SigMod}\) after averaging over the compact
two-plane quotient. The concrete cusp condition for this
\(\mathrm{SL}_2\), including (S1), is supplied by
\srcid{SignatureSplitting.AeFacts.ae\_facts}.
Equation (S6) is its contraction recurrence.
The abstract lemma, through \srcid{SpanPairIff.two\_kill},
therefore gives \(c(hx)=c(x)\).
It follows that \((z-1)C_x(1)=0\), hence \(C_x(1)=0\).

Now fix an originally supplied \(k,f_1\).
Transport its chosen lower splitting by a right translate \(y\).
Applying the same equality at \(\operatorname{lev}(y)k\)
shows that the whole \(\beta\)-component of this splitting
is zero. The chosen plane \(H_1\), complement \(W_2\), and
remaining \(\alpha\)-component witness (S5).
This proves \srcid{TwoStepOf.twoStep\_of'}.
No lower period injectivity, nonzero vector, or
irreducibility argument occurs in this proof.
\end{proof}

The final assembly \srcid{TwoStepOf.twoStep\_of} derives the
two intermediate assertions as follows.
\srcid{SigSplitOf.sigMod\_of} uses \(\operatorname{WeilConv}\)
and \(\operatorname{KDecay}\) to bound the two periods and
their difference.
\srcid{SigSplitOf.sigSplit\_of} uses
\(\operatorname{SplitEx}\), \(\operatorname{PlaneSkew}\),
\(\operatorname{Iwa}\), the closed \(\operatorname{CTBasic}\),
\(\operatorname{IAut}\), and \(\operatorname{SigRe}\).
It first chooses an adapted target model, obtains its
two-component splitting, and transports it back by the
right translation supplied by \(\operatorname{SigRe}\).
Consequently the exact retained inputs of \srcid{twoStep\_of}
are
\[
 \operatorname{WeilConv},\ \operatorname{KDecay},\
 \operatorname{SplitEx},\ \operatorname{PlaneSkew},\
 \operatorname{Iwa},\ \operatorname{IAut},\
 \operatorname{SigRe}.
\]
Their producers are separate parts of the archive.
The stronger upper-vanishing or global-proportionality deductions
are not substituted for the precise conclusion (S5).
